\documentclass[11pt, a4paper]{amsart}
\usepackage{amsmath, amssymb, amsthm, amsfonts}
\usepackage{mathtools}
\usepackage[margin=1in]{geometry}
\usepackage{hyperref}

\newtheorem{theorem}{Theorem}[section]
\newtheorem{lemma}[theorem]{Lemma}
\newtheorem{proposition}[theorem]{Proposition}
\newtheorem{corollary}[theorem]{Corollary}
\newtheorem{fact}[theorem]{Standard Fact}
\theoremstyle{definition}
\newtheorem{definition}[theorem]{Definition}
\newtheorem{remark}[theorem]{Remark}
\newtheorem{notation}[theorem]{Notation}
\newtheorem{hypothesis}[theorem]{Hypothesis}
\newtheorem{axiom}[theorem]{Axiom}

\newcommand{\R}{\mathbb{R}}
\newcommand{\Q}{\mathbb{Q}}
\newcommand{\Z}{\mathbb{Z}}
\newcommand{\HH}{\mathbb{H}}
\newcommand{\disc}{\operatorname{disc}}
\newcommand{\rk}{\operatorname{rank}}
\newcommand{\Gal}{\operatorname{Gal}}
\newcommand{\Frac}{\operatorname{Frac}}
\newcommand{\Hmin}{\mathbf{H}_{\min}}
\newcommand{\Ophi}{\mathbb{Z}[\varphi]}

\title{The 12D Closure Theorem:\\
  Phason Complement Uniqueness and\\
  Self-Reference on the Upward Cascade}
\author{}
\date{}

\begin{document}
\maketitle

\begin{abstract}
Starting from the permeability axiom $\mathbf{F1}$ ($r = 1 + 1/r$,
positive fixed point $\varphi = (1 + \sqrt{5})/2$), we state
three \emph{closure theorems} concerning the upward cascade at
$\Z$-rank $12$, with the following conditionality structure:
C2 (termination, \emph{unconditional} within the admissible
class (A1)--(A4)); C1 (phason-complement $M \cong \Ophi^{2}$,
\emph{conditional on} the minimality hypothesis $\Hmin$
fixing the phason coefficient ring to $\Ophi$); C3
(self-reference, \emph{conditional on} $\Hmin$). Under admissibility (A1)--(A4) together with $\Hmin$, the
$\Z$-rank-$12$ ambient splits as
\[
  L_{12} \;=\; E_8 \oplus M, \qquad M \;\cong\; \Ophi^{2}
\]
as $\Ophi$-modules (so $\rk_\Z M = 4$ and $\rk_\Ophi M = 2$),
and the phason coefficient ring at closure coincides with the
$\mathbf{F1}$ fixed-point ring. The paper uses a single rank convention
throughout: \emph{``$\Z$-rank of $L_{12}$ equals $12$''}
(see Notation~\ref{not:rank-convention}); all statements
involving ``rank'' are over $\Z$ unless explicitly marked
``$\rk_\Ophi$''. The paper has a strict two-layer architecture:

\begin{enumerate}
\item \textbf{Unconditional (Theorem~\ref{thm:C2}, C2).}
  Fix the classical $E_8$ icosian lattice and its
  Elser--Sloane $4+4$ cut-and-project package. Under the
  notion of ``new geometric content'' made precise in
  Definition~\ref{def:no-new}, \emph{no admissible upward
  ambient of $\Z$-rank greater than~$12$ adds new geometric
  content to the upward-cascade closure}. The full admissible
  class is fixed by axioms (A1)--(A4) of
  Definition~\ref{def:admissible}; (A4) supplies the finite
  internal image-rank input that makes the rank-counting
  argument operative.
\item \textbf{Conditional on Hypothesis~\ref{hyp:Hmin}
  (C1, C3).} Under the minimality principle $\Hmin$ fixing
  the phason coefficient ring to $\Ophi$:
  \textbf{C1} (Theorem~\ref{thm:C1}) --- the phason
  complement $M$ with $\rk_\Z M = 4$ is torsion-free over
  $\Ophi$, hence $M \cong \Ophi^{2}$ as $\Ophi$-modules;
  \textbf{C3} (Theorem~\ref{thm:C3}) --- the $\Ophi$ ring
  generated by the phason complement at the 12D closure level
  coincides with the ring generated by the $\mathbf{F1}$ fixed
  point.
\end{enumerate}

$\mathbf{F1}$ is used \emph{only} as a named axiom imported
from
\texttt{papers/cascade-correspondence-foundations/foundations.tex}
\S F1; the foundations manuscript is cited solely for the
axiom statement and \emph{not} as a theorem-grade proof
source. No dependence on the infrastructure papers P1--P7; no
claim about $\alpha$, photons, microtubules, spacetime, or
dynamics is made or implied. Classical inputs
(ring of integers of $\Q(\sqrt{5})$, PID / free-module
structure theorem, Elser--Sloane icosian construction of
$E_8$, Conway--Sloane / Moody--Patera cut-and-project
framework) are cited directly to the classical literature.

This paper is the first entry of the $\alpha$-chain: $\alpha$-2
will treat the phason--Coxeter alignment (A, B); $\alpha$-3 will
assemble $\alpha^{-1} = 137 + \pi/87$. Neither downstream result
is stated or used here.
\end{abstract}

\tableofcontents

%=====================================================================
\section{Introduction and Source Audit}
\label{sec:intro}
%=====================================================================

\subsection{What this paper is}

This is a \emph{parallel} upward-cascade paper: it develops the
12D closure structure $L_{12} = E_8 \oplus M$, independently of
the downward infrastructure papers P1--P7
(cascade-sigma-rationality through cascade-hydrodynamic-projection).
P1--P7 treat the descent $E_8 \to H_4 \to \ldots$; this paper
treats the ascent $E_8 \to L_{12}$. The two sides meet at $E_8$
but do not depend on each other's proof stacks.

The paper takes $\mathbf{F1}$ (the permeability fixed-point axiom
$r = 1 + 1/r$) as its only non-classical input. $\mathbf{F1}$ is
imported as a \emph{named axiom} from
\texttt{papers/cascade-correspondence-foundations/foundations.tex}
\S F1. The foundations manuscript is not itself Clay-bar
rigorous, so we cite it only for the axiom statement and not as
a proof source.

\subsection{What this paper proves}

\textbf{Unconditional layer.}
\begin{description}
\item[\textbf{C2}] (Theorem~\ref{thm:C2}). The upward cascade
  terminates at rank~12. Formally: fix the classical
  $E_8$-icosian lattice with its Elser--Sloane $4+4$
  cut-and-project package
  (Definitions~\ref{def:E8-iso}, \ref{def:cap}); then under
  the notion of ``new geometric content'' made precise in
  Definition~\ref{def:no-new}, no \emph{admissible upward
  ambient} (Def.~\ref{def:admissible}, axioms (A1)--(A4)) of
  $\Z$-rank greater than $12$ adds new geometric content. The
  finite-rank input on the internal image is encoded as
  admissibility axiom (A4); see Remark~\ref{rmk:A4-justification}
  for the discussion of (A4)'s axiomatic status.
\end{description}

\textbf{Conditional on Hypothesis~\ref{hyp:Hmin}:}
\begin{description}
\item[\textbf{C1}] (Theorem~\ref{thm:C1}). Under the minimality
  principle $\Hmin$ fixing the phason coefficient ring to
  $\Ophi$, the phason complement $M$ in $L_{12} = E_8 \oplus M$
  (with $\rk_\Z M = 4$, equivalently $\rk_\Ophi M = 2$) is
  torsion-free over $\Ophi$; by the PID/free-module structure
  theorem (Lemma~\ref{lem:pid-free}), $M \cong \Ophi^{2}$ as
  $\Ophi$-modules.
\item[\textbf{C3}] (Theorem~\ref{thm:C3}). Under the same
  hypothesis, the $\Ophi$ ring generated by the phason
  complement at the 12D closure level coincides with the ring
  generated by the $\mathbf{F1}$ fixed point. This is a
  ring-identification statement; no physical interpretation is
  attached.
\end{description}

\subsection{What this paper does NOT claim}

\begin{enumerate}
\item No unconditional claim that $\Hmin$ holds. The minimality
  principle is stated as a hypothesis; a partial reduction of
  $\Hmin$ to C2 (narrowing the candidate coefficient ring down
  to $\Ophi$ up to the degenerate quartic-order case) is given
  as Proposition~\ref{prop:reduction} at proof-outline level and
  is honestly flagged as conditional.
\item No claim about the fine-structure constant $\alpha$,
  photons, microtubules, or any downstream physical
  application. Those are the subjects of separate papers.
\item No claim that $L_{12}$ is a classified new exceptional
  lattice in the sense of lattice taxonomy; we state and prove
  only $\Z$-module-level / $\Ophi$-module-level statements.
\item No embedding uniqueness of $M$ inside $L_{12}$: we prove
  module-isomorphism-level uniqueness; embedding uniqueness
  would require additional hypotheses not assumed here.
\item No claim that the foundations manuscript is theorem-grade.
  We use $\mathbf{F1}$ only as a named axiom.
\item No Lorentzian, spacetime, gauge-theoretic, or dynamical
  interpretation of the phason complement.
\end{enumerate}

\subsection{Source audit}

\textbf{Axiom-level (cascade).}
$\mathbf{F1}$ from
\texttt{papers/cascade-correspondence-foundations/foundations.tex}
\S F1: the permeability fixed-point axiom $r = 1 + 1/r$ with
positive solution $\varphi = (1+\sqrt{5})/2$. Used as a named
axiom only.

\textbf{Classical references.}
\begin{itemize}
\item \emph{Algebraic number theory.}
  \cite{Marcus, Neukirch} for the ring of integers
  $\mathcal{O}_{\Q(\sqrt{5})} = \Ophi$ identification, the
  discriminant $\disc(\Q(\sqrt{5})) = 5$, the Minkowski bound,
  and the class-number-1 argument giving PID.
\item \emph{Module theory over a PID.}
  \cite{AtiyahMacdonald, DummitFoote} for finitely generated
  torsion-free modules over a PID being free; and the
  structure theorem for finitely generated modules over a PID.
\item \emph{Icosian / $E_8$ / $H_4$ interface.}
  \cite[Ch.~8, Eq.~(17) and Table~8.1]{ConwaySloane},
  \cite[Thm.~4.1, Prop.~4.2]{MoodyPatera}, and
  \cite{ElserSloane} for the classical icosian realization of
  $E_8$ and the $4+4$ cut-and-project package.
\end{itemize}

We do not cite P1--P7, cascade-derivation notes, scripts, or
notebooks as proof sources. The cascade-derivation working notes
are author working notes; their content is converted to clean
theorem statements with classical citations here.

\subsection{Organization}

\S\ref{sec:prelims} states $\mathbf{F1}$ as an axiom, sets up
algebraic number theory for $\Q(\sqrt{5})$ and $\Ophi$, and
proves Lemma~\ref{lem:pid} ($\Ophi$ is a PID),
Lemma~\ref{lem:pid-free} (torsion-free rank-$n$ $\Ophi$-modules
are free).
\S\ref{sec:icosian} records the classical icosian realization
of $E_8$ and the Elser--Sloane $4+4$ cut-and-project package.
\S\ref{sec:C2} proves Theorem~\ref{thm:C2} (C2: termination at
12D) unconditionally.
\S\ref{sec:Hmin} states Hypothesis~\ref{hyp:Hmin} ($\Hmin$) and
includes the conditional Proposition~\ref{prop:reduction}
attempting to reduce $\Hmin$ to C2.
\S\ref{sec:C1} proves Theorem~\ref{thm:C1} (C1) conditional on
$\Hmin$.
\S\ref{sec:C3} proves Theorem~\ref{thm:C3} (C3) conditional on
$\Hmin$.
\S\ref{sec:audit} is the self-audit.
\S\ref{sec:handoff} is the citation contract.
\S\ref{sec:notation} is the notation audit.

%=====================================================================
\section{Preliminaries: $\mathbf{F1}$ and $\Q(\sqrt{5})$ Ring Theory}
\label{sec:prelims}
%=====================================================================

\begin{axiom}[$\mathbf{F1}$: the permeability fixed-point axiom]
\label{ax:F1}
There exists a positive real $\varphi$ satisfying
$\varphi = 1 + 1/\varphi$. This forces $\varphi^2 = \varphi + 1$,
whence
\[
  \varphi \;=\; \frac{1 + \sqrt{5}}{2}.
\]
We import this axiom from
\cite[\S F1]{CascadeFoundations}; the cited manuscript states
the axiom but is not itself Clay-bar rigorous, so we use the
citation \emph{only} to state the axiom.
\end{axiom}

\begin{notation}[The field $K$, ring $\Ophi$, and Galois
  involution $\sigma$]
\label{not:K}
Let $K := \Q(\sqrt{5})$, viewed as a degree-2 number field
over $\Q$. Let $\mathcal{O}_K$ denote its ring of integers and
write $\Ophi := \Z[\varphi]$, the subring of $K$ generated by
$\Z$ and $\varphi$. Let $\sigma \in \Gal(K/\Q)$ be the
nontrivial Galois automorphism, $\sigma(\sqrt{5}) = -\sqrt{5}$,
equivalently $\sigma(\varphi) = 1 - \varphi = -1/\varphi$.
\end{notation}

\begin{notation}[Rank convention]
\label{not:rank-convention}
Throughout this paper, ``rank'' and $\rk$ always mean $\rk_\Z$
(rank as a $\Z$-module), unless marked explicitly $\rk_\Ophi$
(rank as an $\Ophi$-module). In particular, ``$L_{12}$'' and
``$12$D'' refer to $\Z$-rank $12$, not $\Ophi$-rank. Since
$\Ophi$ has $\rk_\Z \Ophi = 2$, an $\Ophi$-module $M$ free of
$\Ophi$-rank $k$ has $\rk_\Z M = 2k$; in particular, an
$\Ophi$-module of $\Z$-rank $4$ has $\Ophi$-rank $2$ (and is
denoted $\Ophi^2$ under this convention, not $\Ophi^4$).
\end{notation}

\begin{lemma}[$\mathcal{O}_K = \Ophi$]
\label{lem:ring-of-integers}
The ring of integers of $K = \Q(\sqrt{5})$ is exactly
$\Z[\varphi]$:
$\mathcal{O}_K = \Z[(1 + \sqrt{5})/2] = \Ophi$.
Equivalently, $\mathcal{O}_K = \{a + b\varphi : a, b \in \Z\}$.
\end{lemma}

\begin{proof}
Classical: for a squarefree integer $d \equiv 1 \pmod{4}$, the
ring of integers of $\Q(\sqrt{d})$ is
$\Z[(1 + \sqrt{d})/2]$. Since $5 \equiv 1 \pmod{4}$, we get
$\mathcal{O}_K = \Z[(1 + \sqrt{5})/2] = \Z[\varphi]$. See
\cite[Ch.~2, Thm.~1]{Marcus} or \cite[Prop.~I.2.4]{Neukirch}.
\end{proof}

\begin{lemma}[$\disc(K) = 5$]
\label{lem:discriminant}
The discriminant of $K = \Q(\sqrt{5})$ is $\disc(K) = 5$.
\end{lemma}

\begin{proof}
For $K = \Q(\sqrt{d})$ with $d \equiv 1 \pmod{4}$ squarefree,
$\disc(K) = d$. See
\cite[Ch.~2, Cor.~to Thm.~1]{Marcus} or
\cite[Ex.~I.2.8]{Neukirch}. With $d = 5$, $\disc(K) = 5$.
\end{proof}

\begin{lemma}[$\Ophi$ is a PID via class-number $1$]
\label{lem:pid}
The ring $\Ophi$ is a principal ideal domain: every nonzero
ideal of $\Ophi$ is principal.
\end{lemma}

\begin{proof}
By Lemma~\ref{lem:ring-of-integers}, $\Ophi = \mathcal{O}_K$,
the ring of integers of $K = \Q(\sqrt{5})$. By
Lemma~\ref{lem:discriminant}, $\disc(K) = 5$.

The Minkowski bound for $K$ (a real quadratic field) is
\[
  M_K \;=\; \sqrt{|\disc(K)|} \bigl(2/\pi\bigr)^{s}
  \;\cdot\; (n!)/n^{n},
\]
where $n = [K : \Q] = 2$ and $s = $ (number of complex embeddings)
$= 0$ (since $K$ is totally real). Thus
\[
  M_K \;=\; \sqrt{5} \cdot \frac{2!}{2^2} \;=\; \frac{\sqrt{5}}{2}
  \;\approx\; 1.118.
\]
Since $M_K < 2$, every ideal class of $\mathcal{O}_K$ contains
an ideal of norm $< 2$; the only such ideals are the unit ideal
$(1) = \mathcal{O}_K$. Hence the class number $h_K = 1$ and
$\mathcal{O}_K$ is a PID. See
\cite[Ch.~5, Thm.~35]{Marcus} for the general Minkowski-bound
argument and
\cite[Ch.~1, \S I.5]{Neukirch} for the discriminant-and-bound
derivation.
\end{proof}

\begin{lemma}[Torsion-free rank-$n$ $\Ophi$-modules are free]
\label{lem:pid-free}
Every finitely generated torsion-free $\Ophi$-module is free.
In particular, every torsion-free rank-$n$ $\Ophi$-module is
isomorphic to $\Ophi^{n}$ as $\Ophi$-modules.
\end{lemma}

\begin{proof}
By Lemma~\ref{lem:pid}, $\Ophi$ is a PID. Over a PID, every
finitely generated torsion-free module is free
\cite[Prop.~3.7 and structure theorem for finitely
generated modules over a PID]{AtiyahMacdonald}; equivalently, by
\cite[Ch.~12, Thm.~4]{DummitFoote}. The rank is the dimension
over the fraction field $K$, so a torsion-free rank-$n$
$\Ophi$-module is isomorphic to $\Ophi^{n}$.
\end{proof}

%=====================================================================
\section{The Classical Icosian / $E_8$ / $H_4$ Interface}
\label{sec:icosian}
%=====================================================================

\begin{definition}[The icosian ring $I$]
\label{def:icosian}
The icosian ring $I$ is the ring of quaternions
\[
  I \;:=\; \Z[\varphi][i, j, k] \;\subset\; \HH,
\]
where $i, j, k$ are the Hamilton quaternion units
($i^2 = j^2 = k^2 = ijk = -1$). As a $\Z$-module, $I$ is free of
rank 8; as an $\Ophi$-module, $I$ is free of rank 4. See
\cite[Ch.~8, Eq.~(17)]{ConwaySloane} for the explicit icosian
realization and the classical identification of the $120$
unit icosians with the binary icosahedral group $2I$.
\end{definition}

\begin{fact}[Classical realization of $E_8$ via icosians]
\label{fact:E8-iso}
The $E_8$ root lattice is realized inside $I \oplus I \subset
\HH \oplus \HH$ as
\[
  E_8 \;=\; \{(x, \sigma(x)) : x \in I\}
  \;\subset\; \HH \oplus \HH \;\cong\; \R^{8},
\]
where $\sigma$ denotes the coefficient-wise Galois action on
$\Ophi$ and hence on $I$. The 240 minimal vectors of $E_8$
correspond to pairs $(\epsilon, \sigma(\epsilon))$ with
$\epsilon \in I$ a unit icosian.
See \cite[Ch.~8, Eq.~(17), Table~8.1]{ConwaySloane} and
\cite[Thm.~4.1]{MoodyPatera}.
\end{fact}

\begin{definition}[$E_8$ in its icosian realization]
\label{def:E8-iso}
$E_8$ denotes the 8-dimensional even unimodular lattice as
realized in Fact~\ref{fact:E8-iso}. We work throughout in this
specific realization.
\end{definition}

\begin{fact}[The Elser--Sloane $4+4$ cut-and-project package]
\label{fact:cap}
Let $\pi_{\rm phys} \colon \HH \oplus \HH \to \HH$ be the
first-factor projection onto the ``physical'' 4-space, and
$\pi_{\rm int} \colon \HH \oplus \HH \to \HH$ the second-factor
projection onto the ``internal'' 4-space. Choose an acceptance
window $W \subset \HH_{\rm int}$ with $\Ophi$-boundary
coordinates (so $W$ has golden-ratio acceptance geometry).
The cut-and-project
\[
  Q_{\rm H_4} \;:=\; \pi_{\rm phys}\bigl(\{v \in E_8 :
  \pi_{\rm int}(v) \in W\}\bigr)
\]
yields the $H_4$ quasicrystal. See \cite{ElserSloane} for the
original construction and
\cite[Thm.~4.1, Prop.~4.2]{MoodyPatera} for the group-theoretic
formulation.
\end{fact}

\begin{definition}[The Elser--Sloane 4+4 cut-and-project package]
\label{def:cap}
The data $(E_8, \HH_{\rm phys}, \HH_{\rm int}, \pi_{\rm phys},
\pi_{\rm int}, W, \sigma)$ of Fact~\ref{fact:cap}. We call this
the \emph{classical package} throughout.
\end{definition}

\begin{remark}[Role of the classical $4+4$ splitting in C2]
\label{rmk:4plus4}
The classical $E_8$ cut-and-project package gives a $4 + 4$
splitting $\HH_{\rm phys} \oplus \HH_{\rm int}$ with both
factors of real dimension $4$ (equivalently, $\rk_\Ophi I =
\rk_\Ophi \sigma(I) = 4$). This fact is classical (Elser--
Sloane; Conway--Sloane \S 8), not a consequence of an
$\Ophi$-module rank calculation internal to this paper.

However, the $12 = 8 + 4$ dimension count in C2 does \emph{not}
follow from the classical $4 + 4$ splitting alone: C2's actual
termination bound comes from admissibility axiom (A4), which
constrains the $\Ophi$-rank of $\pi^{\rm amb}_{\rm int}(M_{\rm
amb})$ to be $\leq 2$. The classical $4 + 4$ data supplies the
concrete ambient-host inside which A4's bound is realised, and
the $\Z$-rank~$2 \cdot 2 = 4$ of the phason-complement image is
where the ``$+4$'' in ``$8 + 4 = 12$'' materialises, but the
logical driver of the termination is (A4), not the $4 + 4$
splitting as a standalone fact. See Theorem~\ref{thm:C2} for
the precise statement.
\end{remark}

%=====================================================================
\section{$\mathbf{C2}$: Termination at 12D (Unconditional)}
\label{sec:C2}
%=====================================================================

\subsection{Making ``new geometric content'' precise}

\begin{definition}[Admissible ambient]
\label{def:admissible}
An \emph{admissible upward ambient} for the classical
Elser--Sloane package (Def.~\ref{def:cap}) is a triple
$(L_{\rm amb}, \mathcal{D}, \Pi^{\rm amb})$ where:
\begin{enumerate}
\item[\textup{(A1)}] $L_{\rm amb}$ is a finitely generated
  torsion-free $\Ophi$-module with $\rk_\Z L_{\rm amb} = N
  \geq 8$. (By Lemma~\ref{lem:pid-free} applied over $\Ophi$,
  $L_{\rm amb}$ is free as an $\Ophi$-module; in particular $N$
  is even.)
\item[\textup{(A2)}] $\mathcal{D}$ is a fixed $\Ophi$-module
  decomposition $L_{\rm amb} = E_8 \oplus M_{\rm amb}$, where
  $E_8$ is the classical icosian-realized $E_8$
  (Def.~\ref{def:E8-iso}) equipped with its inherited
  $\Ophi$-module structure (as a $\Z$-submodule of $I \oplus
  I$ stabilized by the Galois twist). Since $\sigma$ acts
  naturally on $\Ophi$, this action extends entrywise to
  $L_{\rm amb}$, preserving the decomposition and restricting
  to the classical action on $E_8$.
\item[\textup{(A3)}] $\Pi^{\rm amb} = (\pi_{\rm phys}^{\rm amb},
  \pi_{\rm int}^{\rm amb})$ is a pair of $\R$-linear maps
  \[
    \pi_{\rm phys}^{\rm amb} \colon L_{\rm amb} \otimes_\Z \R
      \;\longrightarrow\; \HH_{\rm phys},
    \qquad
    \pi_{\rm int}^{\rm amb} \colon L_{\rm amb} \otimes_\Z \R
      \;\longrightarrow\; \HH_{\rm int},
  \]
  with the following scalar-compatibility structure inherited
  from the classical Elser--Sloane embedding $E_8 = \{(x,
  \sigma(x)) : x \in I\} \subset I \oplus I$ (where $\Ophi$
  acts on $I \oplus I$ diagonally in its natural quaternion
  embedding):
  \begin{enumerate}
  \item[\textup{(A3.a)}] $\pi_{\rm phys}^{\rm amb}$ is
    $\Ophi$-linear:
    $\pi_{\rm phys}^{\rm amb}(\alpha \cdot v) = \alpha \cdot
    \pi_{\rm phys}^{\rm amb}(v)$ for all $\alpha \in \Ophi$,
    $v \in L_{\rm amb}$.
  \item[\textup{(A3.b)}] $\pi_{\rm int}^{\rm amb}$ is
    \emph{$\sigma$-twisted $\Ophi$-linear}:
    $\pi_{\rm int}^{\rm amb}(\alpha \cdot v) = \sigma(\alpha)
    \cdot \pi_{\rm int}^{\rm amb}(v)$ for all $\alpha \in
    \Ophi$, $v \in L_{\rm amb}$.
  \item[\textup{(A3.c)}]
    $\pi_{\rm phys}^{\rm amb}|_{E_8 \otimes \R} = \pi_{\rm phys}$,
    $\pi_{\rm int}^{\rm amb}|_{E_8 \otimes \R} = \pi_{\rm int}$
    (the extensions restrict to the classical package on $E_8$).
  \item[\textup{(A3.d)}] $\pi_{\rm phys}^{\rm amb}|_{M_{\rm amb}
    \otimes \R} = 0$: the phason complement does not contribute
    to the physical target space; its role is only to enlarge
    the acceptance-window / internal parameterization.
  \end{enumerate}
\item[\textup{(A4)}] \emph{(internal image-rank bound)} The
  image $\pi_{\rm int}^{\rm amb}(M_{\rm amb}) \subset \HH_{\rm
  int}$ is a finitely generated $\Ophi$-submodule (where
  $\HH_{\rm int}$ is viewed as an $\Ophi$-module via the
  $\sigma$-twisted action of (A3.b)) of $\Ophi$-rank at most
  $2$, equivalently $\Z$-rank at most $4$.
\end{enumerate}
\end{definition}

\begin{remark}[On (A3.b): the $\sigma$-twist as a compatibility
  axiom]
\label{rmk:sigma-twist}
The $\sigma$-twist in (A3.b) is a \emph{compatibility axiom},
not a derived consequence of the classical Elser--Sloane
package alone. The motivation is the following well-known
algebraic feature of the package: the classical embedding
$E_8 = \{(x, \sigma(x)) : x \in I\} \subset I \oplus I$ is
\emph{not} preserved by the naive diagonal $\Ophi$-action
$\alpha \cdot (x, y) = (\alpha x, \alpha y)$ unless $\alpha
\in \Z$ (since $\sigma(\alpha x) = \sigma(\alpha) \cdot
\sigma(x) \neq \alpha \cdot \sigma(x)$ in general). To make
the embedding $E_8 \hookrightarrow I \oplus I$ equivariant
under a full $\Ophi$-action, one must instead use the
\emph{twisted-diagonal} action $\alpha \cdot (x, y) := (\alpha
x, \sigma(\alpha) y)$. Under this twisted action the
projections become: $\pi_{\rm phys}$ is $\Ophi$-linear and
$\pi_{\rm int}$ is $\sigma$-twisted $\Ophi$-linear, giving
the canonical extension we encode as (A3.a)--(A3.b). See
\cite[Ch.~8, \S2.2]{ConwaySloane} for the icosian form of
this observation and \cite[pp.~6162--6163]{ElserSloane} for
the cut-and-project form. Admissibility \emph{imposes} this
twisted compatibility on the upward extension; we do not
claim it is forced by the bare classical package.
\end{remark}

\begin{remark}[On (A4): why the image-rank bound is part of
  admissibility]
\label{rmk:A4-justification}
(A4) is an explicit constraint, not a derivation: it
encodes the standing interpretation that the upward cascade
extends the classical $4 + 4$ Elser--Sloane package by a
\emph{compatible} phason complement, whose internal image
respects the rank-$2$ $\Ophi$-target dictated by the
acceptance-window structure. Without (A4), an admissible
ambient could in principle have $\pi_{\rm int}^{\rm
amb}(M_{\rm amb})$ of arbitrarily large $\Ophi$-rank inside
$\HH_{\rm int}$ (which is uncountable-rank over $\Ophi$ as
an $\R$-vector space), and the rank-$12$ closure statement
of Theorem~\ref{thm:C2} would be vacuous. (A4) cuts
admissibility down to ambients that can plausibly be
called ``$4 + 4$-compatible'' upward extensions of the
classical package.
\end{remark}

\begin{remark}[On (A3)]
The asymmetry $\pi_{\rm phys}^{\rm amb}|_{M_{\rm amb}} = 0$
formalizes the standing interpretation of the phason
complement $M_{\rm amb}$ in the upward-cascade literature: it
parameterizes the internal (acceptance-window) geometry only,
and does not add new physical-space directions beyond those
already carried by $E_8$. This is the interpretation under
which ``no new geometric content'' has meaning: $E_8$ already
saturates $\HH_{\rm phys} \oplus \HH_{\rm int} = \R^8$ as a
$\Z$-rank-8 lattice, so the only role for additional rank is
to enlarge the acceptance-window parameterization.
\end{remark}

\begin{definition}[Translation-trivial and new geometric content]
\label{def:no-new}
Let $(L_{\rm amb}, \mathcal{D}, \Pi^{\rm amb})$ be admissible
(Def.~\ref{def:admissible}). Define the \emph{extended
internal image} as
\[
  M^{\rm eff} \;:=\; \pi_{\rm int}^{\rm amb}(M_{\rm amb})
    \;\subset\; \HH_{\rm int},
\]
a finitely generated $\Z$-submodule of $\HH_{\rm int} \cong
\R^4$ (not necessarily discrete; in particular, ``sublattice''
in the discrete-subgroup sense is not asserted). We say a
direction $d \in M_{\rm amb}$ is \emph{translation-trivial}
if $\pi_{\rm int}^{\rm amb}(d) = 0$ (equivalently, $d \in
\ker(\pi_{\rm int}^{\rm amb}|_{M_{\rm amb}})$), and we say $d$
\emph{adds new geometric content} otherwise. The subspace of
translation-trivial directions is
\[
  M_{\rm trans} \;:=\; \ker\bigl(\pi_{\rm int}^{\rm amb}
    |_{M_{\rm amb}}\bigr)
    \;\subseteq\; M_{\rm amb}.
\]
Since $\pi_{\rm int}^{\rm amb}$ is $\R$-linear on the
$\Z$-module $M_{\rm amb}$, the kernel is a pure $\Z$-submodule
and $M_{\rm trans}$ is torsion-free as a $\Z$-module.
\end{definition}

\begin{remark}[Justification of the terminology]
\label{rmk:trans-triv-justif}
If $d \in M_{\rm trans}$, then for any
acceptance-window-based cut-and-project
(Def.~\ref{def:cap}), translation of $L_{\rm amb}$ by $d$
leaves both $\pi_{\rm phys}^{\rm amb}$ (zero on $M_{\rm amb}$
by (A3)) and $\pi_{\rm int}^{\rm amb}$ (zero on $d$ by
definition of $M_{\rm trans}$) unchanged on the span of
$E_8$; hence the $H_4$-level quasicrystal $Q_{\rm H_4}$ is
invariant under $d$-translation. This is the geometric
content of ``translation-trivial''.
\end{remark}

\begin{theorem}[$\mathbf{C2}$: termination at 12D]
\label{thm:C2}
Fix the classical Elser--Sloane package of
Definition~\ref{def:cap}. Let $(L_{\rm amb}, \mathcal{D},
\Pi^{\rm amb})$ be an admissible upward ambient
(Def.~\ref{def:admissible}) of $\Z$-rank $N$ (so $N$ is even
and $N \geq 8$). Then:
\begin{enumerate}
\item[\textup{(i)}] The content-bearing part of $M_{\rm amb}$
  has $\Z$-rank at most $4$: letting $M_{\rm trans}$ be the
  translation-trivial sublattice (Def.~\ref{def:no-new}),
  $M_{\rm trans}$ is an $\Ophi$-submodule of $M_{\rm amb}$ and
  $M_{\rm amb} / M_{\rm trans}$ is a finitely generated free
  $\Ophi$-module of $\Ophi$-rank $\leq 2$:
  \[
    \rk_\Z(M_{\rm amb} / M_{\rm trans})
    \;=\; \rk_\Z\bigl(\pi_{\rm int}^{\rm amb}(M_{\rm amb})\bigr)
    \;\leq\; 4,
    \qquad
    \rk_\Ophi(M_{\rm amb} / M_{\rm trans}) \;\leq\; 2.
  \]
\item[\textup{(ii)}] If $N > 12$, then $\rk_\Z M_{\rm trans}
  \geq N - 12 > 0$; in particular $M_{\rm trans} \neq 0$ and
  every excess direction beyond $\Z$-rank $12$ is
  translation-trivial in the sense of
  Def.~\ref{def:no-new}.
\item[\textup{(iii)}] \emph{(Requires $N \geq 12$.)} There
  exists an $\Ophi$-module decomposition $L_{\rm amb} = L_{12}
  \oplus M_{\rm trans}'$, with $L_{12} = E_8 \oplus M_0$ of
  $\Z$-rank $12$ (i.e.\ $\Ophi$-rank $6$), $M_0 \subseteq
  M_{\rm amb}$ an $\Ophi$-submodule of $\Z$-rank $4$
  ($\Ophi$-rank $2$), and $M_{\rm trans}' \subseteq M_{\rm
  trans}$ an $\Ophi$-submodule of $\Z$-rank $N - 12$.
\end{enumerate}
Consequently, no admissible ambient of $\Z$-rank $> 12$ adds
new geometric content (Def.~\ref{def:no-new}) beyond the
$\Z$-rank-$12$ sublattice $L_{12}$.
\end{theorem}

\begin{proof}
Throughout, $M_{\rm amb}$ is a finitely generated torsion-free
$\Ophi$-module (a direct summand of the finitely generated
torsion-free $\Ophi$-module $L_{\rm amb}$ by (A1)--(A2)), hence
free over $\Ophi$ by Lemma~\ref{lem:pid-free}; in particular
$N$ is even and $\rk_\Z M_{\rm amb} = N - 8 = 2\,
\rk_\Ophi M_{\rm amb}$.

\emph{Step 1: (i) --- $M_{\rm trans}$ is an
$\Ophi$-submodule, and image-rank bound.}
By (A3.b), $\pi_{\rm int}^{\rm amb}$ is $\sigma$-twisted
$\Ophi$-linear. Hence if $m \in M_{\rm amb}$ with $\pi_{\rm
int}^{\rm amb}(m) = 0$ and $\alpha \in \Ophi$, then
\[
  \pi_{\rm int}^{\rm amb}(\alpha \cdot m)
  \;=\; \sigma(\alpha) \cdot \pi_{\rm int}^{\rm amb}(m)
  \;=\; \sigma(\alpha) \cdot 0
  \;=\; 0,
\]
so $\alpha \cdot m \in M_{\rm trans}$: $M_{\rm trans}$ is an
$\Ophi$-submodule of $M_{\rm amb}$.

The image-rank bound is precisely the content of admissibility
axiom (A4): $\pi_{\rm int}^{\rm amb}(M_{\rm amb})$ is a
finitely generated $\Ophi$-submodule of $\HH_{\rm int}$
(the latter viewed as an $\Ophi$-module via the
$\sigma$-twisted action of (A3.b)) of $\Ophi$-rank $\leq 2$,
equivalently $\Z$-rank $\leq 4$. Since this image is
torsion-free as an $\Ophi$-module (it is a $\Z$-submodule of
the real vector space $\HH_{\rm int}$, hence has no
$\Z$-torsion; and an $\Ophi$-action on a $\Z$-torsion-free
module has no $\Ophi$-torsion either, by the field-norm
argument used in Theorem~\ref{thm:C1}(b) below), it is free
over the PID $\Ophi$ by Lemma~\ref{lem:pid-free}. The first
isomorphism theorem (in the category of $\Ophi$-modules) gives
\[
  M_{\rm amb} / M_{\rm trans}
  \;\cong\; \pi_{\rm int}^{\rm amb}(M_{\rm amb}),
\]
so $M_{\rm amb} / M_{\rm trans}$ is free of $\Ophi$-rank
$\leq 2$, hence $\Z$-rank $\leq 4$.

\emph{Step 2: (ii) --- the rank bookkeeping.}
Since $M_{\rm amb} / M_{\rm trans}$ is free (hence projective)
as an $\Ophi$-module, the short exact sequence
\[
  0 \;\longrightarrow\; M_{\rm trans}
    \;\longrightarrow\; M_{\rm amb}
    \;\longrightarrow\; M_{\rm amb} / M_{\rm trans}
    \;\longrightarrow\; 0
\]
of $\Ophi$-modules \emph{splits} as $\Ophi$-modules. Hence
\[
  \rk_\Z M_{\rm trans}
  \;=\; \rk_\Z M_{\rm amb} - \rk_\Z(M_{\rm amb} / M_{\rm trans})
  \;\geq\; (N - 8) - 4
  \;=\; N - 12.
\]
For $N > 12$, $\rk_\Z M_{\rm trans} > 0$ and every $d \in
M_{\rm trans}$ is translation-trivial by definition.

\emph{Step 3: (iii) --- explicit $\Ophi$-module splitting,
$N \geq 12$.}
By Step~1, $M_{\rm amb} / M_{\rm trans}$ is free over $\Ophi$
of $\Ophi$-rank $r \in \{0, 1, 2\}$, and the short exact
sequence above splits: pick any $\Ophi$-linear section
$s \colon M_{\rm amb} / M_{\rm trans} \to M_{\rm amb}$ of the
quotient map and set $M_0^{(0)} := s(M_{\rm amb} / M_{\rm
trans})$. Then $M_0^{(0)}$ is an $\Ophi$-submodule of $M_{\rm
amb}$, free of $\Ophi$-rank $r$ and $\Z$-rank $2r \leq 4$, and
$M_{\rm amb} = M_0^{(0)} \oplus M_{\rm trans}$ as
$\Ophi$-modules.

If $r = 2$ (equivalently, $\rk_\Z M_0^{(0)} = 4$), set $M_0 :=
M_0^{(0)}$ and $M_{\rm trans}' := M_{\rm trans}$, and we are
done. If $r < 2$, $M_0^{(0)}$ has $\Z$-rank $< 4$ and we
augment as follows. Since $N \geq 12$, $\rk_\Ophi M_{\rm amb}
= (N-8)/2 \geq 2 \geq r$, so $\rk_\Ophi M_{\rm trans} =
\rk_\Ophi M_{\rm amb} - r \geq 2 - r$. By
Lemma~\ref{lem:pid-free}, $M_{\rm trans}$ is free over $\Ophi$;
choose an $\Ophi$-basis $\{e_1, \ldots, e_{\rk_\Ophi M_{\rm
trans}}\}$ and define
\[
  M_{\rm trans}'' \;:=\; \bigoplus_{i=1}^{2-r} \Ophi \cdot e_i,
  \qquad
  M_{\rm trans}' \;:=\;
  \bigoplus_{i=3-r}^{\rk_\Ophi M_{\rm trans}} \Ophi \cdot e_i.
\]
Then $M_{\rm trans} = M_{\rm trans}'' \oplus M_{\rm trans}'$
tautologically (splitting along a chosen basis), with
$\rk_\Ophi M_{\rm trans}'' = 2 - r$ and $\rk_\Ophi M_{\rm
trans}' = \rk_\Ophi M_{\rm trans} - (2 - r) = (N-8)/2 - 2 =
(N-12)/2$. Set $M_0 := M_0^{(0)} \oplus M_{\rm trans}''
\subseteq M_{\rm amb}$: an $\Ophi$-submodule of $\Ophi$-rank
$r + (2-r) = 2$, $\Z$-rank $4$. Then
\[
  M_{\rm amb}
  \;=\; M_0^{(0)} \oplus M_{\rm trans}'' \oplus M_{\rm trans}'
  \;=\; M_0 \oplus M_{\rm trans}'.
\]
Note $M_{\rm trans}' \subseteq M_{\rm trans}$ consists of
translation-trivial directions. Setting $L_{12} := E_8 \oplus
M_0$ gives $L_{\rm amb} = L_{12} \oplus M_{\rm trans}'$ with
$\rk_\Z L_{12} = 8 + 4 = 12$ as required.

\emph{Step 4: conclusion.}
By Def.~\ref{def:no-new} and Remark~\ref{rmk:trans-triv-justif},
every direction in $M_{\rm trans}'$ is translation-trivial:
adding or subtracting such directions leaves the
$H_4$-level quasicrystal $Q_{\rm H_4}$ unchanged. Hence, for
$N > 12$, no admissible ambient of $\Z$-rank $N$ adds new
geometric content beyond the rank-12 sublattice $L_{12}$.
\end{proof}

\begin{remark}[The sense in which ``12 is maximal'']
\label{rmk:C2-sense}
Theorem~\ref{thm:C2} does \emph{not} claim: ``$L_{12}$ is a
unique exceptional 12-dimensional lattice.'' It claims: once
the classical Elser--Sloane package is fixed, no upward ambient
of rank $> 12$ adds structure relevant to the $E_8 \to H_4$
cut-and-project. This is a termination statement at the level of
``geometric content'' as defined in Definition~\ref{def:no-new},
not a lattice-classification statement. In particular, we
\emph{do not} claim $L_{12}$ is unique up to $\Z$-lattice
isomorphism; uniqueness is established only at the
$\Ophi$-module level in Theorem~\ref{thm:C1} and requires
Hypothesis~\ref{hyp:Hmin}.
\end{remark}

%=====================================================================
\section{The Minimality Principle $\Hmin$}
\label{sec:Hmin}
%=====================================================================

\begin{hypothesis}[$\Hmin$: minimality of the phason coefficient ring]
\label{hyp:Hmin}
Among all rings $R$ satisfying:
\begin{itemize}
\item $\Z \subseteq R \subseteq \mathcal{O}_L$ for some number
  field $L \supseteq \Q$;
\item $\varphi \in R$ (so $\Ophi \subseteq R$);
\item the phason complement $M$ in an admissible ambient
  (Def.~\ref{def:admissible}) carries an $R$-module structure
  compatible with the classical package's $\sigma$-twist
  (Fact~\ref{fact:cap});
\end{itemize}
the \emph{minimal} such $R$ is $\Ophi$ itself.
\end{hypothesis}

\begin{remark}[Status of $\Hmin$]
\label{rmk:Hmin-status}
$\Hmin$ is a structural assumption: it asserts that the phason
complement uses the smallest possible ring-extension of $\Z$
containing $\varphi$. We record a \emph{partial} reduction of
$\Hmin$ to C2 in Proposition~\ref{prop:reduction}: we show that
C2 plus the $\Z$-rank constraint on $M$ already rules out all
candidate rings $R$ of $\Z$-rank $\geq 3$, and (as a classical
fact about quadratic orders) rules out all $\Z$-rank-2
candidates containing $\varphi$ except $\Ophi$. A single
residual gap remains: the $\rk_\Z R = 4$ case (possible quartic
orders containing $\varphi$), which is not excluded by $\Z$-rank
alone; that remaining case is flagged explicitly in
Remark~\ref{rmk:posture}. We do not claim the reduction is
Clay-bar-complete; Theorems~\ref{thm:C1}, \ref{thm:C3} are
stated \emph{conditional on $\Hmin$} throughout.
\end{remark}

\begin{proposition}[Partial reduction of $\Hmin$ to C2 under
  (R1), proof outline]
\label{prop:reduction}
Fix an admissible ambient $(L_{\rm amb}, \mathcal{D},
\Pi^{\rm amb})$ of $\Z$-rank exactly $12$ (so $\rk_\Z M_{\rm
amb} = 4$; we make no claim about $M_{\rm trans}$ vanishing
in this case), and let $R$ be a ring satisfying the
hypotheses of $\Hmin$ (i.e.\
$\Z \subseteq R \subseteq \mathcal{O}_L$ for some number field
$L \supseteq \Q$, $\varphi \in R$, $M_{\rm amb}$ carries a
compatible $R$-module structure). Assume additionally:
\begin{itemize}
\item[\textup{(R1)}] $M_{\rm amb}$ is \emph{free} as an
  $R$-module. (Equivalently: torsion-free of rank $k_R :=
  \rk_R M_{\rm amb}$ over the PID $R$; this requires $R$ to be
  a PID, which holds for all orders in imaginary quadratic
  fields of class number $1$ and for $R = \Ophi$ by
  Lemma~\ref{lem:pid}.)
\end{itemize}
Then $k_R \geq 1$ and the $\Z$-rank identity
\[
  \rk_\Z M_{\rm amb}
  \;=\; k_R \cdot \rk_\Z R
  \;=\; 4
\]
forces $\rk_\Z R \in \{1, 2, 4\}$. Separating cases:
\begin{enumerate}
\item[\textup{(a)}] $\rk_\Z R = 1$. Then $R = \Z$, which does
  not contain $\varphi$, contradicting the $\Hmin$-hypothesis
  $\varphi \in R$. Ruled out.
\item[\textup{(b)}] $\rk_\Z R = 2$. Then $\Frac(R)$ is a
  quadratic extension of $\Q$ containing $\varphi$, hence
  $\Frac(R) = \Q(\varphi) = K$. Then $R \subseteq \mathcal{O}_K
  = \Ophi$ (Lemma~\ref{lem:ring-of-integers}); conversely
  $\Ophi \subseteq R$ since $\varphi \in R$ and $1 \in R$. Hence
  $R = \Ophi$.
\item[\textup{(c)}] $\rk_\Z R = 4$. Not excluded by the
  $\Z$-rank count alone; $R$ could be an order of $\Z$-rank $4$
  in a degree-$4$ number field containing $\varphi$, in which
  case $k_R = 1$ and $M_{\rm amb} \cong R$ as $R$-modules. This
  is the residual case the proposition does not exclude.
\end{enumerate}
Thus under (R1), any $R$ satisfying $\Hmin$'s hypotheses must
have $\rk_\Z R \in \{2, 4\}$, and in the $\rk_\Z R = 2$ case
$R = \Ophi$ is forced.
\end{proposition}

\begin{proof}[Proof outline]
(R1) is explicit. The $\Z$-rank identity
$\rk_\Z M_{\rm amb} = k_R \cdot \rk_\Z R$ holds for any free
$R$-module of rank $k_R$. Since $\rk_\Z M_{\rm amb} = 4$ (the
standing assumption that $L_{\rm amb}$ has $\Z$-rank exactly
$12$, with $\rk_\Z E_8 = 8$), $\rk_\Z R$ divides $4$, giving
$\rk_\Z R \in \{1, 2, 4\}$.

Case (a) is immediate. Case (b) uses classical quadratic
number theory: the only order in $\Q(\varphi) = \Q(\sqrt 5)$
containing the ring of integers $\Ophi = \mathcal O_K$
(Lemma~\ref{lem:ring-of-integers}) is $\Ophi$ itself, and any
subring of $\Ophi$ containing $\varphi$ is $\Ophi$. Case (c)
requires additional structure beyond $\Z$-rank and is not
pursued here.
\end{proof}

\begin{remark}[Residual gap and posture of the paper]
\label{rmk:posture}
Proposition~\ref{prop:reduction} does \emph{not} establish
$\Hmin$ unconditionally: the quartic-order case
$\rk_\Z R = 4$ is not ruled out by rank counting. A natural
exclusion argument would use that in case (c), $k_R = 1$ and
$M_{\rm amb} \cong R$ has a distinguished multiplicative
structure (from $R$'s ring structure) not supported by the
classical Elser--Sloane package; but formalizing this
requires a \emph{further} axiom about compatibility with the
cut-and-project's $\sigma$-twist structure, which we decline
to introduce here. Accordingly, we state Theorems
\ref{thm:C1}, \ref{thm:C3} \emph{conditional on $\Hmin$},
viewing Proposition~\ref{prop:reduction} as a pointer to
where the reduction can be closed in a downstream paper, not
as closure itself.
\end{remark}

%=====================================================================
\section{$\mathbf{C1}$: Phason Complement Module Uniqueness (Conditional)}
\label{sec:C1}
%=====================================================================

\begin{theorem}[$\mathbf{C1}$: phason complement is $\Ophi^2$ under $\Hmin$]
\label{thm:C1}
Assume Hypothesis~\ref{hyp:Hmin} ($\Hmin$). Let $(L_{\rm amb},
\mathcal{D}, \Pi^{\rm amb})$ be a 12D admissible ambient
(Def.~\ref{def:admissible}, $\rk_\Z L_{\rm amb} = 12$) with
phason complement $M := M_{\rm amb}$ (so $\rk_\Z M = 4$). Then:
\begin{enumerate}
\item[\textup{(a)}] $M$ carries a canonical $\Ophi$-module
  structure inherited from the admissibility datum (A1)--(A2),
  with $\rk_\Ophi M = 2$ (equivalently, $\rk_\Z M = 2 \cdot
  \rk_\Z \Ophi = 4$ under Notation~\ref{not:rank-convention}).
\item[\textup{(b)}] $M$ is torsion-free as an $\Ophi$-module.
\item[\textup{(c)}] Consequently, $M \cong \Ophi^{2}$ as an
  $\Ophi$-module.
\end{enumerate}
\end{theorem}

\begin{proof}
\emph{(a):} By admissibility (A1), $M$ is a finitely generated
torsion-free $\Ophi$-module. The $\Z$-rank identity
$\rk_\Z M = \rk_\Ophi M \cdot \rk_\Z \Ophi$ (for free
$\Ophi$-modules; follows from $\Ophi$ free of $\Z$-rank $2$ and
Lemma~\ref{lem:pid-free}) combined with $\rk_\Z M = 4$
(Theorem~\ref{thm:C2}(i) with $N = 12$, or the
$12 = \rk_\Z E_8 + \rk_\Z M = 8 + 4$ split) and $\rk_\Z \Ophi =
2$ (Lemma~\ref{lem:ring-of-integers}) gives $\rk_\Ophi M = 2$.
The hypothesis $\Hmin$ ensures that the \emph{coefficient ring}
(i.e., the ring over which $M$ is viewed as a module) is
$\Ophi$ and not a larger ring; this is what makes the
$\Ophi$-rank a well-defined invariant here.

\emph{(b) (norm argument):} We must show that if $\alpha m = 0$
for $\alpha \in \Ophi$ and $m \in M$, then either $\alpha = 0$
or $m = 0$. Suppose $\alpha m = 0$ with $\alpha \neq 0$. Apply
$\sigma(\alpha) \in \Ophi$ (where $\sigma \in \Gal(K/\Q)$ is
the Galois involution of Notation~\ref{not:K}):
\[
  \sigma(\alpha) \cdot (\alpha m) \;=\; N_{K/\Q}(\alpha) \cdot m
  \;=\; 0,
\]
where $N_{K/\Q}(\alpha) := \alpha \cdot \sigma(\alpha) \in \Z
\setminus \{0\}$ (nonzero because $\Ophi$ is a domain,
$\alpha \neq 0$, and $\sigma(\alpha) \neq 0$; see
\cite[Ch.~2, \S2.3]{Marcus}). This gives $n \cdot m = 0$ for the
integer $n := N_{K/\Q}(\alpha) \neq 0$, contradicting $M$ being
torsion-free as a $\Z$-module (which follows from (A1): $M$ is a
$\Z$-direct-summand of the torsion-free $\Z$-module $L_{\rm amb}$).
Hence $m = 0$, and $M$ is torsion-free over $\Ophi$.

\emph{(c):} By Lemma~\ref{lem:pid-free} applied with $n =
\rk_\Ophi M = 2$, $M \cong \Ophi^{2}$ as $\Ophi$-module.
\end{proof}

\begin{remark}[The specific role of $\Hmin$ in Theorem~\ref{thm:C1}]
\label{rmk:C1-Hmin-role}
Items (a)--(c) of Theorem~\ref{thm:C1} can be proved from
admissibility (A1)--(A2) alone \emph{given} that the coefficient
ring is $\Ophi$. The work of $\Hmin$ is upstream of (a)--(c):
it ensures that $\Ophi$, rather than some larger ring $R
\supsetneq \Ophi$, is the coefficient ring over which $M$ is
viewed as a module. Without $\Hmin$, the same rank-bookkeeping
+ PID-freeness argument would still produce ``$M \cong R^{r}$''
for whatever the actual coefficient ring $R$ turns out to be;
the conditional content ``$R = \Ophi$, so $M \cong \Ophi^2$'' is
the ring-minimality assertion. Accordingly, C1's \emph{freeness}
conclusion is unconditional on $\Hmin$; its \emph{identification
of the ring as $\Ophi$} is the conditional part.
\end{remark}

\begin{remark}[Notation clarification: rank convention]
\label{rmk:notation}
This paper's notation is fixed by
Notation~\ref{not:rank-convention}: ``rank'' means $\rk_\Z$ unless
marked otherwise, so $L_{12}$ has $\Z$-rank~$12$ and $M \cong
\Ophi^{2}$ means $\rk_\Ophi M = 2$ (equivalently $\rk_\Z M = 4$).
With this convention, $L_{12} = E_8 \oplus M_{12}$ has
$\Z$-rank~$12 = 8 + 4$. We emphasise that $\Ophi^{2}$ and
$\Ophi^{4}$ are \emph{distinct} $\Ophi$-modules with different
$\Ophi$-ranks and different $\Z$-ranks; they are \emph{not} the
same underlying algebraic object. Any downstream paper using a
different notation (e.g.\ writing ``$\Ophi^{4}$'' for the same
phason complement in a different rank convention) must open with
an explicit conversion statement identifying $\rk_\Ophi$ and
$\rk_\Z$ in both conventions. This paper commits to
$M_{12} \cong \Ophi^{2}$ as $\Ophi$-module, $\rk_\Z M_{12} = 4$.
\end{remark}

%=====================================================================
\section{$\mathbf{C3}$: Self-Reference Closure (Conditional)}
\label{sec:C3}
%=====================================================================

\begin{theorem}[$\mathbf{C3}$: self-reference closure under $\Hmin$]
\label{thm:C3}
Assume Hypothesis~\ref{hyp:Hmin} ($\Hmin$). Then:
\begin{enumerate}
\item[\textup{(a)}] The coefficient ring of the phason
  complement $M$ at the 12D closure level is $\Ophi$.
\item[\textup{(b)}] The ring generated by the permeability
  fixed point $\varphi$ of Axiom~\ref{ax:F1} (i.e., the
  smallest subring of $\R$ containing $\Z$ and $\varphi$) is
  $\Ophi$.
\item[\textup{(c)}] The two rings coincide: both are $\Ophi$.
\end{enumerate}
This is a ring-identification theorem. No physical
interpretation, no dynamics, no downstream $\alpha$/photon/
microtubule claim is attached.
\end{theorem}

\begin{proof}
\emph{(a):} Theorem~\ref{thm:C1} under $\Hmin$ identifies $M
\cong \Ophi^{2}$ as an $\Ophi$-module. The minimality of the
coefficient ring — i.e.\ that $\Ophi$ is the coefficient ring,
not merely \emph{a} compatible ring — is the content of
$\Hmin$ itself (Hypothesis~\ref{hyp:Hmin}), which is assumed.
Thus (a) reduces to $\Hmin$ together with C1's freeness result
(cf.\ Remark~\ref{rmk:C1-Hmin-role} for the separation of the
two roles of $\Hmin$ and (A1)).

\emph{(b):} $\varphi$ satisfies $\varphi^2 = \varphi + 1$, so
$\Z[\varphi] = \Z + \Z\varphi = \Ophi$ (the last equality by
Lemma~\ref{lem:ring-of-integers}, since $\Ophi$ is the ring
of integers of $\Q(\sqrt{5})$). The smallest subring of $\R$
containing $\Z$ and $\varphi$ is precisely $\Z[\varphi] = \Ophi$.

\emph{(c):} Both rings identified in (a) and (b) equal $\Ophi$.
\end{proof}

\begin{remark}[``Permeability describes permeability'']
\label{rmk:self-ref}
A colloquial way of stating C3 is: ``Permeability describes
permeability'' --- the axiom's fixed-point ring reappears as
the coefficient ring at the 12D closure level. This is an
\emph{interpretation} of Theorem~\ref{thm:C3}, not part of the
theorem statement. The theorem itself is the ring identification.
Any metaphysical or physical reading of this interpretation is
outside the scope of this paper.
\end{remark}

%=====================================================================
\section{Self-Audit}
\label{sec:audit}
%=====================================================================

\textbf{Unconditional layer (Theorem~\ref{thm:C2}):}
\begin{itemize}
\item[[1]] C2 proved by literal rank-counting within the
  admissible class fixed by Def.~\ref{def:admissible}, axioms
  (A1)--(A4); (A4) is the finite internal-image-rank
  hypothesis that makes the rank count operative
  (Remark~\ref{rmk:A4-justification}). Proof in
  \S\ref{sec:C2}.
\item[[2]] ``No new geometric content'' made precise in
  Def.~\ref{def:no-new}; not left as slogan.
\item[[3]] 4 + 4 splitting of Elser--Sloane package cited from
  \cite{ConwaySloane, MoodyPatera, ElserSloane}.
\item[[4]] $\Z[\varphi]$ = PID proved from Minkowski bound
  (Lemma~\ref{lem:pid}), not asserted.
\item[[5]] PID/free-module structure theorem cited with
  specific reference (Lemma~\ref{lem:pid-free}).
\end{itemize}

\textbf{Conditional on $\Hmin$ (Theorems~\ref{thm:C1}, \ref{thm:C3}):}
\begin{itemize}
\item[[6]] $\Hmin$ stated explicitly as
  Hypothesis~\ref{hyp:Hmin}; not hidden.
\item[[7]] Proposition~\ref{prop:reduction} gives a \emph{partial}
  reduction of $\Hmin$ to C2 + (R1): the candidate rings are
  cut down from all number-field orders containing $\varphi$
  to just $\Ophi$ and possible $\Z$-rank-$4$ quartic orders.
  The residual $\rk_\Z R = 4$ case is explicitly flagged
  (Remark~\ref{rmk:posture}); the ``proof outline'' label is
  honest about what is and is not proved.
\item[[8]] C1 conditional fencing explicit in theorem
  statement; uniqueness stated only at $\Ophi$-module level
  (not embedding level). $\Ophi$-torsion-freeness proved via
  the field-norm argument (Thm.~\ref{thm:C1}(b)), not
  tacitly asserted from ``$\Ophi$ has no zero-divisors''.
\item[[9]] C3 conditional on $\Hmin$; ring-identification only;
  no physical interpretation attached.
\item[[10]] Rank convention is fixed once in
  Notation~\ref{not:rank-convention} and used uniformly
  throughout: all ranks are $\Z$-ranks unless marked
  $\rk_\Ophi$. The informal cascade-derivation notes'
  ``$\Ophi^{4}$'' labelling is flagged in
  Remark~\ref{rmk:notation} as a distinct (equivalent) reading
  that belongs to the $\alpha$-2 paper.
\end{itemize}

\textbf{Not claimed:}
\begin{itemize}
\item No unconditional C1 or C3.
\item No claim $\Hmin$ is proved.
\item No $\alpha$, photon, microtubule, spacetime, dynamics.
\item No dependence on P1--P7.
\item No lattice-classification uniqueness of $L_{12}$.
\item No embedding uniqueness of $M$ in $L_{12}$.
\end{itemize}

%=====================================================================
\section{Citation Contract and Downstream Handoff}
\label{sec:handoff}
%=====================================================================

\textbf{Unconditionally citable:}
\begin{itemize}
\item Theorem~\ref{thm:C2} (C2: termination at 12D).
\item Lemma~\ref{lem:pid} ($\Ophi$ PID).
\item Lemma~\ref{lem:pid-free} (torsion-free rank-$n$
  $\Ophi$-modules are free).
\item Classical package (Def.~\ref{def:cap},
  Fact~\ref{fact:E8-iso}) --- imported from cited sources.
\end{itemize}

\textbf{Conditional on $\Hmin$ (NOT Clay-bar unconditionally
citable):}
\begin{itemize}
\item Theorem~\ref{thm:C1} (C1: $M \cong \Ophi^{2}$ as
  $\Ophi$-module).
\item Theorem~\ref{thm:C3} (C3: self-reference closure).
\end{itemize}

\textbf{Downstream usage ($\alpha$-2, $\alpha$-3):} Downstream
papers may cite C2 unconditionally; C1 and C3 only conditionally,
with the explicit $\Hmin$ dependence noted. The $\alpha$-2 paper
(phason-Coxeter alignment) reconciles the $\Z$-rank vs
$\Ophi$-rank convention in its own preamble; we do not attempt
the reconciliation beyond Remark~\ref{rmk:notation} here.

\textbf{Does NOT close:} the minimality principle as an
unconditional theorem. Proposition~\ref{prop:reduction} closes
the $\rk_\Z R \in \{1, 2\}$ cases under hypothesis (R1): case
$\rk_\Z R = 1$ is ruled out (no $\varphi$), and case $\rk_\Z R
= 2$ forces $R = \Ophi$ by quadratic number theory. The single
residual case is $\rk_\Z R = 4$ (quartic orders containing
$\varphi$), deferred to a follow-up paper as noted in
Remark~\ref{rmk:posture}.

%=====================================================================
\section{Notation Audit}
\label{sec:notation}
%=====================================================================

\begin{center}
\begin{tabular}{|l|l|l|}
\hline
Symbol & Meaning & First use \\
\hline
$\varphi$ & $(1+\sqrt 5)/2$, positive fixed point of
$\mathbf{F1}$ & Axiom~\ref{ax:F1} \\
$K$ & $\Q(\sqrt 5)$ & Notation~\ref{not:K} \\
$\mathcal O_K$ & ring of integers of $K$ & Notation~\ref{not:K} \\
$\Ophi$ & $\Z[\varphi] = \mathcal O_K$ & Notation~\ref{not:K} \\
$\sigma$ & Galois involution of $K/\Q$; $\varphi \mapsto 1-\varphi$ &
  Notation~\ref{not:K} \\
$\disc(K)$ & discriminant of $K$ & Lemma~\ref{lem:discriminant} \\
$h_K$ & class number of $K$ & Lemma~\ref{lem:pid} \\
$I$ & icosian ring & Definition~\ref{def:icosian} \\
$E_8$ & $E_8$ lattice in its icosian realization &
  Definition~\ref{def:E8-iso} \\
$\HH$ & Hamilton quaternions & Definition~\ref{def:icosian} \\
$\pi_{\rm phys}, \pi_{\rm int}$ & physical / internal projections
  in the classical package & Fact~\ref{fact:cap} \\
$W$ & acceptance window in the cut-and-project & Fact~\ref{fact:cap} \\
$Q_{\rm H_4}$ & $H_4$ quasicrystal from Elser--Sloane &
  Fact~\ref{fact:cap} \\
$L_{\rm amb}, M_{\rm amb}$ & admissible upward ambient and its phason
  complement & Def.~\ref{def:admissible} \\
$L_{12}, M$ & 12D closure ambient and its (rank-4 $\Z$) phason
  complement & Theorem~\ref{thm:C2} \\
$M_{\rm trans}$ & translation-trivial excess above rank 12 &
  Theorem~\ref{thm:C2} \\
$\Hmin$ & minimality hypothesis on the phason coefficient ring &
  Hypothesis~\ref{hyp:Hmin} \\
C1, C2, C3 & Theorems~\ref{thm:C1}, \ref{thm:C2}, \ref{thm:C3}
  respectively & \\
\hline
\end{tabular}
\end{center}

\begin{thebibliography}{99}

\bibitem{CascadeFoundations}
\emph{Cascade Correspondence Foundations},
\texttt{papers/cascade-correspondence-foundations/foundations.tex}.
Cited solely for the statement of Axiom $\mathbf{F1}$ (\S F1).
\emph{Not} cited as a theorem-grade proof source.

\bibitem{Marcus}
D.~A.~Marcus, \emph{Number Fields}, 2nd ed., Universitext,
Springer, 2018. Ring of integers of $\Q(\sqrt d)$ for $d \equiv
1 \pmod 4$: Ch.~2, Thm.~1 (and its Corollary for the
discriminant formula). Minkowski bound and class-number
finiteness: Ch.~5, Thm.~35. Field norm and its integrality:
Ch.~2, \S2.3.

\bibitem{Neukirch}
J.~Neukirch, \emph{Algebraic Number Theory}, Grundlehren der
mathematischen Wissenschaften 322, Springer, 1999. Ring of
integers of quadratic fields: Prop.~I.2.4 and Ex.~I.2.8.
Minkowski / discriminant bound for the class number:
\S I.5 (in particular the bound used in the proof of
Theorem~I.5.3, the finiteness of the class group).

\bibitem{AtiyahMacdonald}
M.~F.~Atiyah, I.~G.~Macdonald, \emph{Introduction to Commutative
Algebra}, Addison-Wesley, 1969. Finitely generated torsion-free
modules over a PID are free: combine Prop.~3.7 (localization of
torsion) with Ch.~8 exercises on modules over PIDs; see also
\cite{DummitFoote} below for the complete structure theorem.

\bibitem{DummitFoote}
D.~S.~Dummit, R.~M.~Foote, \emph{Abstract Algebra}, 3rd ed.,
Wiley, 2004. Fundamental Theorem for finitely generated modules
over a PID (including the torsion-free $\Rightarrow$ free
corollary): Ch.~12, \S12.1, Theorems 4 and 5 (invariant-factor
and elementary-divisor forms).

\bibitem{ConwaySloane}
J.~H.~Conway, N.~J.~A.~Sloane, \emph{Sphere Packings,
Lattices and Groups}, 3rd ed., Springer, 1999. Icosian ring
$I \subset \HH$ and the unit icosians realizing $2I$: Ch.~8,
\S2.1 (``The icosian ring''), with the explicit $\Z[\varphi]$
generator list displayed on p.~207. $E_8 \subset \HH \oplus
\HH$ via icosians: Ch.~8, \S2.2 and Eq.~(8.2.17)
(``$E_8$ as icosians''); Table~8.1 (p.~209) enumerates the
$240$ minimal vectors.

\bibitem{MoodyPatera}
R.~V.~Moody, J.~Patera, \emph{Quasicrystals and icosians},
J.~Phys.~A \textbf{26} (1993), 2829--2853.
The explicit $E_8 \to H_4$ cut-and-project interface: Thm.~4.1
(structure of the projection onto the physical $H_4$-root
system) and Prop.~4.2 (acceptance-window description of the
resulting quasicrystal).

\bibitem{ElserSloane}
V.~Elser, N.~J.~A.~Sloane, \emph{A highly symmetric
four-dimensional quasicrystal}, J.~Phys.~A \textbf{20} (1987),
6161--6167. Original cut-and-project construction of the
$H_4$-symmetric 4D quasicrystal from $E_8$: \S2 (pp.~6162--6163)
sets up the $I \oplus \sigma(I)$ decomposition of $E_8$;
\S3 (pp.~6163--6165) defines the acceptance window and projects
onto $\HH_{\rm phys}$; the main diagram of projections is given
in Fig.~1, p.~6162.

\end{thebibliography}

\end{document}
